\section{Lista de endpoints}
\label{anexo:listaendpoints}

Las siguientes tablas presentan los \textit{endpoints} implementados en el \textit{Backend}.

\small
\renewcommand{\arraystretch}{1.15}

\newcommand{\epmethod}[1]{\texttt{#1}}
\newcommand{\eproute}[1]{\path{#1}}

\subsection*{Autenticación}
\begin{table}[H]
	\centering
	\begin{tabular}{|p{1.8cm}|p{4.4cm}|p{7.6cm}|}
		\hline
		\textbf{Método} &   \textbf{Ruta}                   &   \textbf{Función}                                                    \\ \hline
		\epmethod{POST} &   \eproute{/auth/register}        &   Registrar un nuevo usuario en la plataforma.                        \\ \hline
		\epmethod{POST} &   \eproute{/auth/login}           &   Autenticar al usuario e iniciar sesión.                             \\ \hline
		\epmethod{POST} &   \eproute{/auth/refresh}         &   Renovar el \textit{access token} usando el \textit{refresh token}.  \\ \hline
		\epmethod{POST} &   \eproute{/auth/verify-email}    &   Confirmar la verificación del correo electrónico.                   \\ \hline
		\epmethod{POST} &   \eproute{/auth/forgot-password} &   Solicitar el envío de token para restablecer contraseña.            \\ \hline
		\epmethod{POST} &   \eproute{/auth/reset-password}  &   Restablecer la contraseña usando el token de recuperación.          \\ \hline
		\epmethod{POST} &   \eproute{/auth/logout}          &   Cerrar una sesión específica del usuario.                           \\ \hline
		\epmethod{POST} &   \eproute{/auth/logout-all}      &   Cerrar todas las sesiones activas del usuario.                      \\ \hline
		\epmethod{GET}  &   \eproute{/auth/sessions}        &   Listar las sesiones activas del usuario autenticado.                \\ \hline
	\end{tabular}
	\caption[Endpoints de Autenticación]{Endpoints del módulo de autenticación.}
	\label{tab:endpointsBackend}
\end{table}

\subsection*{Usuarios}
\begin{table}[H]
	\centering
	\begin{tabular}{|p{1.8cm}|p{4.4cm}|p{7.6cm}|}
		\hline
		\textbf{Método}     & \textbf{Ruta}                         & \textbf{Función}                                          \\ \hline
		\epmethod{GET}      & \eproute{/users/}                     & Listar estudiantes disponibles en la plataforma.          \\ \hline
		\epmethod{GET}      & \eproute{/users/me}                   & Obtener la información del perfil del usuario autenticado.\\ \hline
		\epmethod{GET}      & \eproute{/users/:id/profile}          & Obtener el perfil público de un usuario por su ID.        \\ \hline
		\epmethod{PUT}      & \eproute{/users/:id}                  & Actualizar la información de un usuario.                  \\ \hline
		\epmethod{PUT}      & \eproute{/users/:id/change-password}  & Cambiar la contraseña del usuario.                        \\ \hline
		\epmethod{DELETE}   & \eproute{/users/:id}                  & Eliminar la cuenta de usuario.                            \\ \hline
		\epmethod{PUT}      & \eproute{/users/:id/profile-picture}  & Actualizar la fotografía de perfil del usuario.           \\ \hline
		\epmethod{GET}      & \eproute{/pending-organizations}      & Listar organizaciones externas pendientes de aprobación.  \\ \hline
		\epmethod{PATCH}    & \eproute{/users/:id/approve}          & Aprobar o rechazar cuentas de organizaciones externas.    \\ \hline
	\end{tabular}
	\caption[Endpoints de Usuarios]{Endpoints del módulo de usuarios.}
\end{table}

\subsection*{Proyectos}
\begin{table}[H]
	\centering
	\begin{tabular}{|p{1.8cm}|p{4.4cm}|p{7.6cm}|}
		\hline
		\textbf{Método}     & \textbf{Ruta}                                     & \textbf{Función}                                      \\ \hline
		\epmethod{GET}      & \eproute{/projects}                               & Listar proyectos aplicando filtros de búsqueda.       \\ \hline
		\epmethod{GET}      & \eproute{/projects/my-projects}                   & Listar proyectos creados por el usuario autenticado.  \\ \hline
		\epmethod{GET}      & \eproute{/projects/pending-approvals}             & Listar proyectos pendientes de aprobación o rechazo.  \\ \hline
		\epmethod{POST}     & \eproute{/projects/:projectId/approve}            & Aprobar o rechazar un proyecto específico.            \\ \hline
		\epmethod{GET}      & \eproute{/projects/:projectId}                    & Obtener el detalle de un proyecto por ID.             \\ \hline
		\epmethod{POST}     & \eproute{/projects}                               & Crear un nuevo proyecto.                              \\ \hline
		\epmethod{PUT}      & \eproute{/projects/:projectId}                    & Actualizar la información de un proyecto existente.   \\ \hline
		\epmethod{POST}     & \eproute{/projects/:projectId/resources}          & Agregar archivos de recurso a un proyecto.            \\ \hline
		\epmethod{DELETE}   & \eproute{/projects/:projectId/resources/:fileId}  & Eliminar un recurso específico de un proyecto.        \\ \hline
		\epmethod{PATCH}    & \eproute{/projects/:projectId/status}             & Cambiar el estado de un proyecto.                     \\ \hline
		\epmethod{DELETE}   & \eproute{/projects/:projectId}                    & Eliminar un proyecto.                                 \\ \hline
	\end{tabular}
	\caption[Endpoints de Proyectos]{Endpoints del módulo de proyectos.}
\end{table}

\subsection*{Aplicaciones}
\begin{table}[H]
	\centering
	\begin{tabular}{|p{1.8cm}|p{4.4cm}|p{7.6cm}|}
		\hline
		\textbf{Método}     & \textbf{Ruta}                                                     & \textbf{Función}                                              \\ \hline
		\epmethod{POST}     & \eproute{/applications}                                           & Registrar una postulación de estudiante a un proyecto.        \\ \hline
		\epmethod{GET}      & \eproute{/applications/completed}                                 & Obtener el historial de proyectos finalizados del estudiante. \\ \hline
		\epmethod{GET}      & \eproute{/applications/project/:projectId}                        & Obtener las postulaciones de un proyecto específico.          \\ \hline
		\epmethod{GET}      & \eproute{/applications/my-current-project}                        & Obtener el proyecto actual del estudiante autenticado.        \\ \hline
		\epmethod{GET}      & \eproute{/applications/project/:projectId/participants}           & Listar participantes aceptados en un proyecto.                \\ \hline
		\epmethod{PATCH}    & \eproute{/applications/:applicationId/status}                     & Aprobar o rechazar una postulación.                           \\ \hline
		\epmethod{PATCH}    & \eproute{/applications/project/:projectId/leave}                  & Permitir que un estudiante abandone un proyecto.              \\ \hline
		\epmethod{PATCH}    & \eproute{/applications/project/:projectId/participant/:studentId} & Eliminar a un participante de un proyecto.                    \\ \hline
	\end{tabular}
	\caption[Endpoints de Aplicaciones]{Endpoints del módulo de aplicaciones.}
\end{table}

\subsection*{Horas}
\begin{table}[H]
	\centering
	\begin{tabular}{|p{1.8cm}|p{4.4cm}|p{7.6cm}|}
		\hline
		\textbf{Método}     & \textbf{Ruta}                                 & \textbf{Función}                                              \\ \hline
		\epmethod{POST}     & \eproute{/hours}                              & Registrar horas trabajadas con sus evidencias.                \\ \hline
		\epmethod{GET}      & \eproute{/hours/total}                        & Obtener el total de horas aprobadas del estudiante.           \\ \hline
		\epmethod{GET}      & \eproute{/hours/:hourId}                      & Obtener el detalle de un registro de horas.                   \\ \hline
		\epmethod{GET}      & \eproute{/hours/project/:projectId/pending}   & Listar horas pendientes de aprobación en un proyecto.         \\ \hline
		\epmethod{PATCH}    & \eproute{/hours/:hourId/status}               & Aprobar o rechazar un registro de horas.                      \\ \hline
		\epmethod{POST}     & \eproute{/hours/:hourId/evidences}            & Agregar evidencias a un registro de horas.                    \\ \hline
		\epmethod{DELETE}   & \eproute{/hours/:hourId/evidences/:fileId}    & Eliminar una evidencia específica de un registro de horas.    \\ \hline
		\epmethod{PATCH}    & \eproute{/hours/:hourId/resubmit}             & Reenviar horas rechazadas para nueva revisión.                \\ \hline
	\end{tabular}
	\caption[Endpoints de Horas]{Endpoints del módulo de horas.}
\end{table}

\subsection*{Archivos}
\begin{table}[H]
	\centering
	\begin{tabular}{|p{1.8cm}|p{4.4cm}|p{7.6cm}|}
		\hline
		\textbf{Método} & \textbf{Ruta}                         & \textbf{Función}                                                  \\ \hline
		\epmethod{GET}  & \eproute{/files/:id/download}         & Descargar un archivo almacenado en la plataforma.                 \\ \hline
		\epmethod{GET}  & \eproute{/files/:id/preview}          & Previsualizar un archivo sin descargarlo.                         \\ \hline
		\epmethod{POST} & \eproute{/files/editor/upload-temp}   & Subir un archivo temporal para el editor de texto enriquecido.    \\ \hline
	\end{tabular}
	\caption[Endpoints de Archivos]{Endpoints del módulo de archivos.}
\end{table}

\subsection*{Otros}
\begin{table}[H]
	\centering
	\begin{tabular}{|p{1.8cm}|p{4.4cm}|p{7.6cm}|}
		\hline
		\textbf{Método} & \textbf{Ruta}         & \textbf{Función}              \\ \hline
		\epmethod{GET}  & \eproute{/faculties}  & Listar facultades disponibles.\\ \hline
		\epmethod{GET}  & \eproute{/careers}    & Listar carreras disponibles.  \\ \hline
		\epmethod{GET}  & \eproute{/campuses}   & Listar campus disponibles.    \\ \hline
		\epmethod{GET}  & \eproute{/tags}       & Listar etiquetas disponibles. \\ \hline
	\end{tabular}
	\caption[Endpoints complementarios]{Endpoints complementarios de consulta general.}
\end{table}

\normalsize

\section{Código Fuente}

El código referente al \textbf{\textit{Backend}} del proyecto, puede ser obtenido
a través de:
\href{https://github.com/csuvg/control-horas-extension-backend.git}{\textbf{ActiUVG - Backend}}

El código referente al \textbf{\textit{Frontend}} del proyecto, puede ser obtenido
a través de:
\href{https://github.com/csuvg/control-horas-extension-frontend.git}{\textbf{ActiUVG - Frontend}}

\textbf{Nota:} Por motivos de seguridad, ambos repositorios se encuentran privados,
en caso de requerir acceso al código fuente, por favor contactar al autor de ambos repositorios.